%----------------------------------------------------------------------------------------
% EERSTEGRAADSONGELIJKHEDEN, MET INTERVALLEN. Alle opgaven zijn eigen werk.
%
% Opbouw: Probleem, Start, leerstof, en op een nieuwe pagina de Opdrachten en Test jezelf
% (2de graad: geen Koppel). Labels van dit hoofdstuk beginnen met h2-.
% Pakket 4 en 5: ook dubbele ongelijkheden (\alleenvoor{4,5}).
%----------------------------------------------------------------------------------------

\mychapter{2}{Ongelijkheden en intervallen}

\begin{probleem}
\naast{Fitnesszaal Sterk vraagt 25 euro per maand en 2 euro per bezoek. Fitnesszaal Fit vraagt geen maandbedrag, maar 5 euro per bezoek. Bij weinig bezoeken is Fit goedkoper, bij veel bezoeken Sterk. Vanaf hoeveel bezoeken per maand is Sterk goedkoper?}{\small\renewcommand{\arraystretch}{1.3}%
\begin{tabular}{|l|c|c|c|}
\hline
bezoeken & 4 & 8 & 12 \\ \hline
Sterk & 33 & 41 & 49 \\ \hline
Fit & 20 & 40 & 60 \\ \hline
\end{tabular}}
\end{probleem}

\lkr{Leerplan \pakketnaam: LPD \alleenvoor{1,2,3}{15}\alleenvoor{4,5}{19} (eerstegraadsongelijkheden oplossen, de oplossing als interval noteren en voorstellen op de getallenas). Voorkennis: hoofdstuk 1 en het ordenen van getallen op de getallenas. De intervallen staan in de Vlaamse notatie, met de haken naar buiten bij een open grens, en op de getallenas met dezelfde haken. Gebruikt uw school bolletjes (vol of open), vermeld dat dan in de les. \alleenvoor{4,5}{In dit pakket komen ook dubbele ongelijkheden aan bod.}}

In hoofdstuk 1 zochten we het getal dat een vergelijking waar maakt. Maar vaak vraag je niet wanneer twee dingen gelijk zijn, wel wanneer het ene kleiner is dan het andere: vanaf wanneer is een abonnement goedkoper, hoeveel dozen mogen in de lift? Het antwoord is dan geen enkel getal, maar een hele reeks getallen. In dit hoofdstuk leer je die reeks noteren als een interval, en ongelijkheden oplossen met bijna dezelfde regels als vergelijkingen. Bijna, want er is één verrassing.

\begin{oefeningen}[Start]
\oef{h2-start-1}

\oef{h2-verken}
\end{oefeningen}

%========================================================================================
\section{Ongelijkheden en intervallen}

\subsection{Meer dan één oplossing}

In de Verken zocht je welke getallen de ongelijkheid $2x + 1 < 7$ waar maken. Voor $x = 2$ is $2 \cdot 2 + 1 = 5$, en $5 < 7$ klopt. Voor $x = 3$ is het linkerlid 7, en $7 < 7$ klopt niet. Ook $2{,}9$ werkt, en $-100$, en $0{,}5$: elk getal kleiner dan 3. Een ongelijkheid heeft dus meestal niet één oplossing, maar oneindig veel.

\begin{theorie}
Een \textbf{\textcolor{accent}{ongelijkheid}} vergelijkt twee uitdrukkingen met een van de tekens $<$, $\leqslant$, $>$ of $\geqslant$, en er staat een onbekende in. Een \textbf{\textcolor{accent}{oplossing}} is een getal dat, ingevuld voor de onbekende, een ware uitspraak geeft. Alle oplossingen samen vormen de oplossingsverzameling $V$.
\end{theorie}

Je leest $a \leqslant b$ als \textit{$a$ is kleiner dan of gelijk aan $b$}. Dus $7 \leqslant 7$ klopt, maar $7 < 7$ niet. Op de getallenas betekent $a < b$ dat $a$ links van $b$ ligt. Daarom is $-5 < -2$, ook al is 5 groter dan 2.

\voorbeeld{}{Is 4 een oplossing van $3x - 5 \geqslant 2x - 1$? En 3?\label{vb:h2-controle}}{%
\uitleg{We vullen elk getal in beide leden in, en kijken of de uitspraak klopt. Dit voorbeeld draait om het verschil tussen $>$ en $\geqslant$.}
\begin{stappen}
x = 4: \ 3 \cdot 4 - 5 = 7 \ \text{en} \ 2 \cdot 4 - 1 = 7 & $7 \geqslant 7$ klopt: 4 is een oplossing. Bij $>$ was het geen oplossing geweest. \\
x = 3: \ 3 \cdot 3 - 5 = 4 \ \text{en} \ 2 \cdot 3 - 1 = 5 & $4 \geqslant 5$ klopt niet: 3 is geen oplossing.
\end{stappen}}

\begin{oefeningen}
\oef{h2-controle}
\end{oefeningen}

\subsection{Intervallen}

Hoe schrijf je \textit{alle getallen kleiner dan 3} op? Opsommen kan niet. Je kunt het wel tekenen: op de getallenas maak je het stuk links van 3 dik. Het getal 3 zelf hoort er niet bij. Dat toon je met een haak die van het dikke stuk wegkijkt.

\figuren{\begin{tikzpicture}[x=.6cm,y=.6cm]
  \getallenas{-3}{5}{-3,-2,-1,0,1,2,3,4,5}
  \ivdeel{-3.4}{3}\haakl{3}
\end{tikzpicture}}

Zo'n doorlopend stuk van de getallenas heet een interval. Er zijn twee soorten: met twee grenzen, of met één grens.

\begin{theorie}
Neem twee reële getallen $a < b$. Een \textbf{\textcolor{accent}{interval}} met grenzen $a$ en $b$ bevat alle getallen tussen $a$ en $b$. Een grens die erbij hoort, krijgt een haak naar het interval toe. Een grens die er niet bij hoort, krijgt een haak die ervan wegkijkt.

\vspace{1mm}
{\centering\arrayrulecolor{accent}\renewcommand{\arraystretch}{1.5}%
\begin{tabular}{|>{\centering\arraybackslash}m{1.6cm}|>{\centering\arraybackslash}m{2.9cm}|>{\centering\arraybackslash}m{2.1cm}|>{\centering\arraybackslash}m{1.6cm}|>{\centering\arraybackslash}m{3.5cm}|}
\hline
\textbf{notatie} & \textbf{betekenis} & \textbf{naam} & \multicolumn{2}{c|}{\textbf{voorbeeld}} \\ \hline
$\formule{\ivgg ab}$ & $\formule{a \leqslant x \leqslant b}$ & \textbf{\textcolor{accent}{gesloten}} & $\ivgg{-1}{3}$ & \raisebox{-2.8mm}[4.2mm][3.6mm]{\begin{tikzpicture}[x=.38cm,y=.38cm]
  \getallenas{-2}{4}{-1,3}
  \ivdeel{-1}{3}\haakl{-1}\haakr{3}
\end{tikzpicture}} \\ \hline
$\formule{\ivoo ab}$ & $\formule{a < x < b}$ & \textbf{\textcolor{accent}{open}} & $\ivoo{-1}{3}$ & \raisebox{-2.8mm}[4.2mm][3.6mm]{\begin{tikzpicture}[x=.38cm,y=.38cm]
  \getallenas{-2}{4}{-1,3}
  \ivdeel{-1}{3}\haakr{-1}\haakl{3}
\end{tikzpicture}} \\ \hline
$\formule{\ivgo ab}$ & $\formule{a \leqslant x < b}$ & \textbf{\textcolor{accent}{halfopen}} & $\ivgo{-1}{3}$ & \raisebox{-2.8mm}[4.2mm][3.6mm]{\begin{tikzpicture}[x=.38cm,y=.38cm]
  \getallenas{-2}{4}{-1,3}
  \ivdeel{-1}{3}\haakl{-1}\haakl{3}
\end{tikzpicture}} \\ \hline
$\formule{\ivog ab}$ & $\formule{a < x \leqslant b}$ & \textbf{\textcolor{accent}{halfopen}} & $\ivog{-1}{3}$ & \raisebox{-2.8mm}[4.2mm][3.6mm]{\begin{tikzpicture}[x=.38cm,y=.38cm]
  \getallenas{-2}{4}{-1,3}
  \ivdeel{-1}{3}\haakr{-1}\haakr{3}
\end{tikzpicture}} \\ \hline
\end{tabular}\par}
\end{theorie}

De haken in de notatie zijn dezelfde als op de getallenas. Bij $\ivgo{-1}{3}$ hoort $-1$ erbij en 3 niet. Het interval bevat dus $-1$, $0$, $2{,}9$ en $2{,}999$, maar niet 3.

\voorbeeld{}{Schrijf als interval en teken op de getallenas: $-2 < x \leqslant 4$.\label{vb:h2-interval}}{%
\uitleg{De grenzen zijn $-2$ en $4$. Het teken $<$ bij $-2$ zegt dat $-2$ er niet bij hoort, het teken $\leqslant$ bij 4 dat 4 er wel bij hoort.}
\naast{$V = \ivog{-2}{4}$: links een haak die wegkijkt, rechts een haak naar het interval toe.}{\begin{tikzpicture}[x=.55cm,y=.55cm]
  \getallenas{-3}{5}{-3,-2,-1,0,1,2,3,4,5}
  \ivdeel{-2}{4}\haakr{-2}\haakr{4}
\end{tikzpicture}}}

En de getallen kleiner dan 3? Dat interval heeft maar één grens: naar links loopt het door zonder einde. Daarvoor gebruiken we het teken $\infty$, \textit{oneindig}.

\begin{theorie}
Een interval met één grens $a$ loopt naar één kant door tot in het oneindige:

\vspace{1mm}
{\centering\arrayrulecolor{accent}\renewcommand{\arraystretch}{1.5}%
\begin{tabular}{|>{\centering\arraybackslash}m{2.6cm}|>{\centering\arraybackslash}m{3cm}|>{\centering\arraybackslash}m{2.2cm}|>{\centering\arraybackslash}m{3.5cm}|}
\hline
\textbf{notatie} & \textbf{betekenis} & \multicolumn{2}{c|}{\textbf{voorbeeld}} \\ \hline
$\formule{\ivgo a\pinf}$ & $\formule{x \geqslant a}$ & $\ivgo{1}{\pinf}$ & \raisebox{-2.8mm}[4.2mm][3.6mm]{\begin{tikzpicture}[x=.38cm,y=.38cm]
  \getallenas{-2}{4}{0,1}
  \ivdeel{1}{4.5}\haakl{1}
\end{tikzpicture}} \\ \hline
$\formule{\ivoo a\pinf}$ & $\formule{x > a}$ & $\ivoo{1}{\pinf}$ & \raisebox{-2.8mm}[4.2mm][3.6mm]{\begin{tikzpicture}[x=.38cm,y=.38cm]
  \getallenas{-2}{4}{0,1}
  \ivdeel{1}{4.5}\haakr{1}
\end{tikzpicture}} \\ \hline
$\formule{\ivog\minf a}$ & $\formule{x \leqslant a}$ & $\ivog{\minf}{1}$ & \raisebox{-2.8mm}[4.2mm][3.6mm]{\begin{tikzpicture}[x=.38cm,y=.38cm]
  \getallenas{-2}{4}{0,1}
  \ivdeel{-2.4}{1}\haakr{1}
\end{tikzpicture}} \\ \hline
$\formule{\ivoo\minf a}$ & $\formule{x < a}$ & $\ivoo{\minf}{1}$ & \raisebox{-2.8mm}[4.2mm][3.6mm]{\begin{tikzpicture}[x=.38cm,y=.38cm]
  \getallenas{-2}{4}{0,1}
  \ivdeel{-2.4}{1}\haakl{1}
\end{tikzpicture}} \\ \hline
\end{tabular}\par}

\vspace{1mm}
Alle reële getallen samen: $\formule{\R = \ivoo\minf\pinf}$.
\end{theorie}

Oneindig is geen getal: je kunt het nooit bereiken. Daarom staat er bij $\pinf$ en $\minf$ altijd een haak die wegkijkt. De getallen kleiner dan 3 vormen dus het interval $\ivoo{\minf}{3}$.

\voorbeeld{}{Schrijf als interval en teken: $x \geqslant -1$.\label{vb:h2-interval2}}{%
\uitleg{De enige grens is $-1$, en die hoort erbij. Naar rechts loopt het interval door.}
\naast{$V = \ivgo{-1}{\pinf}$: bij $-1$ een haak naar het interval toe, bij $\pinf$ een haak die wegkijkt, want oneindig is geen getal.}{\begin{tikzpicture}[x=.55cm,y=.55cm]
  \getallenas{-3}{5}{-3,-2,-1,0,1,2,3,4,5}
  \ivdeel{-1}{5.4}\haakl{-1}
\end{tikzpicture}}}

\nujij{Schrijf als interval: $x < 0$, en $\tfrac12 \leqslant x \leqslant 2$.}{$\ivoo{\minf}{0}$ en $\ivgg{\tfrac12}{2}$.}

\begin{oefeningen}
\oef{h2-interval}

\oef{h2-aflezen}
\end{oefeningen}

%========================================================================================
\section{Ongelijkheden oplossen}

Hoe vind je zonder gokken alle oplossingen van $2x + 1 < 7$? Zoals bij vergelijkingen: maak de ongelijkheid stap voor stap eenvoudiger, tot er $x < \ldots$ staat. Maar mag dat met dezelfde regels? We gaan ze één voor één na.

\subsection{Bij beide leden optellen}

Neem een ware ongelijkheid, zoals $2 < 5$. Tel bij beide leden 4 op: $6 < 9$, en dat klopt nog altijd. Trek 10 af: $-8 < -5$, en ook dat klopt. Op de getallenas schuiven beide getallen evenveel op, dus het linkse blijft links.

\figuren{\begin{tikzpicture}[x=.55cm,y=.55cm]
  \getallenas{-1}{10}{0,2,5,6,9}
  \draw[->,grijs,thick] (2,.25) to[bend left=35] node[above] {\scriptsize $+4$} (6,.25);
  \draw[->,grijs,thick] (5,.25) to[bend left=35] node[above] {\scriptsize $+4$} (9,.25);
  \fill[accent] (2,0) circle (3pt) (5,0) circle (3pt);
  \fill[grijs] (6,0) circle (3pt) (9,0) circle (3pt);
\end{tikzpicture}}

\begin{theorie}
\begin{stelling}{Bij beide leden optellen}
Tel je bij beide leden van een ongelijkheid hetzelfde getal of dezelfde term op, of trek je het van beide leden af, dan blijft het ongelijkheidsteken staan.

\vspace{1mm}
\qquad $\formule{A < B \iff A + c < B + c}$
\end{stelling}
Hetzelfde geldt voor $\leqslant$, $>$ en $\geqslant$.
\end{theorie}

Ook bij ongelijkheden mag je dus een term overbrengen naar het andere lid, met het tegengestelde teken.

\voorbeeld{}{Los op: $3x + 2 < 2x - 1$.\label{vb:h2-optellen}}{%
\uitleg{We brengen $2x$ naar links en $2$ naar rechts, net als bij een vergelijking.}
\naast{\begin{stappenvrij}{3.6cm}{4.6cm}
3x - 2x < -1 - 2 & Breng over, met het tegengestelde teken. \\
x < -3 & Reken uit. \\
V = \ivoo{\minf}{-3} & Controle met $x = -4$: $-10 < -9$ klopt.
\end{stappenvrij}}{\begin{tikzpicture}[x=.5cm,y=.5cm]
  \getallenas{-6}{1}{-6,-5,-4,-3,-2,-1,0,1}
  \ivdeel{-6.4}{-3}\haakl{-3}
\end{tikzpicture}}}

\subsection{Vermenigvuldigen met een positief getal}

Vermenigvuldig beide leden van $2 < 5$ met 3: $6 < 15$. Deel ze door 2: $1 < 2{,}5$. Het teken blijft kloppen. Vermenigvuldigen met een positief getal rekt de getallenas uit of drukt ze samen, maar wat links lag, blijft links.

\begin{theorie}
\begin{stelling}{Vermenigvuldigen met een positief getal}
Vermenigvuldig of deel je beide leden van een ongelijkheid door hetzelfde positieve getal, dan blijft het ongelijkheidsteken staan.

\vspace{1mm}
\qquad $\formule{A < B \iff A \cdot c < B \cdot c}$ \qquad als $c > 0$
\end{stelling}
\end{theorie}

\voorbeeld{}{Los op: $3x - 5 > 7$.\label{vb:h2-positief}}{%
\uitleg{Eerst brengen we $-5$ over. Daarna delen we door 3, een positief getal: het teken blijft.}
\naast{\begin{stappenvrij}{3.6cm}{4.6cm}
3x > 12 & Breng $-5$ over. \\
x > 4 & Deel door 3. \\
V = \ivoo{4}{\pinf} & Controle met $x = 5$: $10 > 7$.
\end{stappenvrij}}{\begin{tikzpicture}[x=.5cm,y=.5cm]
  \getallenas{1}{8}{1,2,3,4,5,6,7,8}
  \ivdeel{4}{8.4}\haakr{4}
\end{tikzpicture}}}

\begin{oefeningen}
\oef{h2-positief}
\end{oefeningen}

\subsection{Vermenigvuldigen met een negatief getal}

In de Verken vermenigvuldigde je $2 < 5$ met $-1$. Dat geeft $-2$ en $-5$, en nu is $-2 > -5$: het teken klopt niet meer, het is omgekeerd. Waarom? Vermenigvuldigen met $-1$ spiegelt elk getal rond 0. Wat links lag, komt rechts te liggen.

\figuren{\begin{tikzpicture}[x=.55cm,y=.55cm]
  \getallenas{-6}{6}{-5,-2,0,2,5}
  \fill[accent] (2,0) circle (3pt) (5,0) circle (3pt);
  \fill[grijs] (-2,0) circle (3pt) (-5,0) circle (3pt);
  \draw[->,grijs,thick] (2,.3) to[bend right=50] (-2,.3);
  \draw[->,grijs,thick] (5,.3) to[bend right=40] (-5,.3);
  \node[above,font=\scriptsize] at (0,1.9) {$\cdot(-1)$};
\end{tikzpicture}}

Met elk ander negatief getal gebeurt hetzelfde: $2 < 5$ wordt na vermenigvuldigen met $-3$ gelijk aan $-6 > -15$. Dat is de verrassing van dit hoofdstuk.

\begin{theorie}
\begin{stelling}{Vermenigvuldigen met een negatief getal}
Vermenigvuldig of deel je beide leden van een ongelijkheid door hetzelfde negatieve getal, dan keert het ongelijkheidsteken om.

\vspace{1mm}
\qquad $\formule{A < B \iff A \cdot c > B \cdot c}$ \qquad als $c < 0$
\end{stelling}
Zo wordt $\leqslant$ ook $\geqslant$, en omgekeerd.
\end{theorie}

\voorbeeld{}{Los op: $-2x \geqslant 6$.\label{vb:h2-negatief}}{%
\uitleg{We delen door $-2$, een negatief getal: het teken keert om.}
\naast{\begin{stappenvrij}{3.6cm}{4.6cm}
x \leqslant \dfrac{6}{-2} & Deel door $-2$: $\geqslant$ wordt $\leqslant$. \\[2ex]
x \leqslant -3 & Reken uit. \\
V = \ivog{\minf}{-3} & Controle met $x = -4$: $8 \geqslant 6$ klopt. Met $x = 0$: $0 \geqslant 6$ klopt niet.
\end{stappenvrij}}{\begin{tikzpicture}[x=.5cm,y=.5cm]
  \getallenas{-6}{1}{-6,-5,-4,-3,-2,-1,0,1}
  \ivdeel{-6.4}{-3}\haakr{-3}
\end{tikzpicture}}}

Vergeet je het teken om te keren, dan vind je precies de verkeerde getallen. Je kunt het delen door een negatief getal ook vermijden: breng de term met $x$ naar het lid waar ze positief wordt.

\voorbeeld{}{Los op: $5 - 2x \geqslant 11$, op twee manieren.\label{vb:h2-tweemanieren}}{%
\uitleg{Eerst met de termen met $x$ links, en dus delen door $-2$. Daarna met $x$ rechts.}
\begin{stappen}
-2x \geqslant 6 \ \Rightarrow \ x \leqslant -3 & Breng 5 over en deel door $-2$: het teken keert om. \\
5 - 11 \geqslant 2x \ \Rightarrow \ -6 \geqslant 2x & Of: breng $-2x$ naar rechts en 11 naar links. \\
-3 \geqslant x & Deel door 2, positief: het teken blijft. Dat is hetzelfde als $x \leqslant -3$.
\end{stappen}
Beide manieren geven $V = \ivog{\minf}{-3}$.}

\nujij{Los op: $4 - 3x < 10$.}{$-3x < 6$, dus $x > -2$: $V = \ivoo{-2}{\pinf}$.}

\begin{oefeningen}
\oef{h2-negatief}
\end{oefeningen}

\subsection{Een ongelijkheid herleiden}

Nu kun je elke eerstegraadsongelijkheid oplossen. Haakjes en noemers werk je weg zoals bij vergelijkingen: vermenigvuldigen met het kgv mag, want dat is een positief getal.

\voorbeeld{}{Los op: $\dfrac{x - 1}{2} - \dfrac{2x + 3}{3} < 1$.\label{vb:h2-herleid}}{%
\uitleg{Het kgv van de noemers is 6, positief: het teken blijft. Op het einde delen we door een negatief getal.}
\begin{stappen}
3(x - 1) - 2(2x + 3) < 6 & Vermenigvuldig elke term met 6, met haakjes rond de tellers. \\
3x - 3 - 4x - 6 < 6 & Werk de haakjes uit. \\
-x < 15 & Neem samen en breng $-9$ over. \\
x > -15 & Deel door $-1$: het teken keert om. \\
V = \ivoo{-15}{\pinf} & Controle met $x = 0$: $-\tfrac12 - 1 = -1{,}5 < 1$ klopt.
\end{stappen}}

Controleer een ongelijkheid met een getal in $V$ en een getal erbuiten. Het grensgetal zelf geeft een gelijkheid: met $x = -15$ is het linkerlid $-8 + 9 = 1$. Dat bevestigt dat je de grens goed hebt.

\begin{werkwijze}
Werk de noemers en de haakjes weg. Breng de termen met $x$ naar één lid en de getallen naar het andere. Deel door de coëfficiënt van $x$, en keer het teken om als die negatief is. Schrijf $V$ als interval en teken het. Controleer met een getal in $V$ en een getal erbuiten.
\end{werkwijze}

Valt $x$ weg, dan blijft er een uitspraak over die altijd of nooit klopt, zoals bij vergelijkingen.

\voorbeeld{}{Los op: $2(x + 1) > 2x + 5$.\label{vb:h2-geen}}{%
\uitleg{We werken uit en kijken wat er overblijft.}
\begin{stappen}
2x + 2 > 2x + 5 & Werk de haakjes uit. \\
0x > 3 & Breng over: $0 > 3$ klopt voor geen enkele $x$. \\
V = \emptyset &
\end{stappen}}

Bij $3x - 1 \leqslant 3(x + 1)$ krijg je daarentegen $0x \leqslant 4$, en $0 \leqslant 4$ klopt altijd: $V = \R$.

\nujij{Los op: $\dfrac{x}{2} - 1 \geqslant \dfrac{x + 4}{3}$.}{Maal 6: $3x - 6 \geqslant 2x + 8$, dus $x \geqslant 14$: $V = \ivgo{14}{\pinf}$.}

\begin{oefeningen}
\oef{h2-herleid}
\end{oefeningen}

\begin{alleenvoor}{4,5}
\subsection{Dubbele ongelijkheden}

Soms moet een uitdrukking tussen twee grenzen liggen, zoals in $-3 < 2x + 1 \leqslant 7$. Dat zijn twee ongelijkheden tegelijk. Je kunt ze apart oplossen, maar het kan sneller: pas dezelfde bewerking toe op alle drie de delen.

\voorbeeld{}{Los op: $-3 < 2x + 1 \leqslant 7$.\label{vb:h2-dubbel}}{%
\uitleg{We trekken 1 af van de drie delen, en delen ze daarna door 2.}
\naast{\begin{stappenvrij}{3.6cm}{4.6cm}
-4 < 2x \leqslant 6 & Trek 1 af van elk deel. \\
-2 < x \leqslant 3 & Deel elk deel door 2. \\
V = \ivog{-2}{3} &
\end{stappenvrij}}{\begin{tikzpicture}[x=.5cm,y=.5cm]
  \getallenas{-3}{4}{-3,-2,-1,0,1,2,3,4}
  \ivdeel{-2}{3}\haakr{-2}\haakr{3}
\end{tikzpicture}}}

Deel je door een negatief getal, dan keren beide tekens om. Schrijf het resultaat daarna om, zodat het kleinste getal links staat.

\voorbeeld{}{Los op: $1 \leqslant 5 - 2x < 9$.\label{vb:h2-dubbel2}}{%
\uitleg{Na het aftrekken van 5 delen we door $-2$: beide tekens keren om.}
\begin{stappen}
-4 \leqslant -2x < 4 & Trek 5 af van elk deel. \\
2 \geqslant x > -2 & Deel door $-2$: beide tekens keren om. \\
-2 < x \leqslant 2 & Lees van rechts naar links. $V = \ivog{-2}{2}$.
\end{stappen}}

\begin{oefeningen}
\oef{h2-dubbel}
\end{oefeningen}
\end{alleenvoor}

%========================================================================================
\section{Vraagstukken met ongelijkheden}

Wanneer stelt een vraagstuk een ongelijkheid op, en geen vergelijking? Als het vraagt naar \textit{minstens}, \textit{hoogstens}, \textit{vanaf wanneer} of \textit{goedkoper dan}. Het antwoord is dan een reeks waarden.

\subsection{Van tekst naar ongelijkheid}

Een paar woorden komen altijd terug. Let vooral op \textit{minstens} en \textit{hoogstens}: daar hoort de grens zelf bij.

\begin{center}\arrayrulecolor{lijngrijs}\renewcommand{\arraystretch}{1.5}%
\begin{tabular}{|>{\centering\arraybackslash}m{2.6cm}|>{\centering\arraybackslash}m{1.4cm}|>{\centering\arraybackslash}m{7.4cm}|}
\hline
\textbf{in woorden} & \textbf{teken} & \textbf{voorbeeld} \\ \hline
minstens & $\geqslant$ & je moet minstens 1,40~m groot zijn: $l \geqslant 1{,}40$ \\ \hline
hoogstens & $\leqslant$ & een lift draagt hoogstens 630~kg: $g \leqslant 630$ \\ \hline
meer dan & $>$ & een tent voor meer dan 4 personen: $n > 4$ \\ \hline
minder dan & $<$ & het kost minder dan 20 euro: $p < 20$ \\ \hline
niet meer dan & $\leqslant$ & niet meer dan 50~km/h: $v \leqslant 50$ \\ \hline
\end{tabular}\end{center}

\nujij{Schrijf als ongelijkheid: een koffer weegt niet meer dan 23~kg.}{$g \leqslant 23$.}

\subsection{Een vraagstuk oplossen}

Bij \textit{goedkoper dan} vergelijk je twee bedragen: het ene is kleiner dan het andere.

\voorbeeld{}{Terug naar het probleem. Fitnesszaal Sterk vraagt 25 euro per maand en 2 euro per bezoek. Fitnesszaal Fit vraagt geen maandbedrag, maar 5 euro per bezoek. Bij weinig bezoeken is Fit goedkoper, bij veel bezoeken Sterk. Vanaf hoeveel bezoeken per maand is Sterk goedkoper?\label{vb:h2-terug}}{%
\uitleg{We kiezen het aantal bezoeken als onbekende, en schrijven beide bedragen in $x$. Op het einde kijken we welke getallen in deze context kunnen: een aantal bezoeken is een natuurlijk getal.}
\naast{\begin{stappenvrij}{4.1cm}{3.6cm}
x = \text{het aantal bezoeken} & De onbekende. \\
25 + 2x < 5x & Sterk kost minder dan Fit. \\
25 < 3x & Breng $2x$ naar rechts. \\
x > \dfrac{25}{3} \approx 8{,}3 & Deel door 3, positief.
\end{stappenvrij}}{\begin{tikzpicture}[x=.36cm,y=.065cm,font=\scriptsize]
  \draw[black!14,very thin] (0,0) grid[xstep=1,ystep=10] (12,60);
  \draw[->,grijs] (0,0) -- (12.8,0) node[right,text=black] {$x$};
  \draw[->,grijs] (0,0) -- (0,64) node[above,text=black] {euro};
  \foreach \t in {0,4,8,12}{\node[below,inner sep=1.5pt] at (\t,0) {$\t$};}
  \foreach \t in {20,40,60}{\node[left,inner sep=1.5pt] at (0,\t) {$\t$};}
  \foreach \n in {0,...,12}{%
    \fill[accent] (\n,{25+2*\n}) circle (2.2pt);
    \fill[grijs] (\n,{5*\n}) +(-2pt,-2pt) rectangle +(2pt,2pt);}
  \node[accent,right] at (12.2,49) {Sterk};
  \node[grijs,right] at (12.2,58) {Fit};
\end{tikzpicture}}
Een aantal bezoeken is een natuurlijk getal, dus $x \geqslant 9$. \textit{Vanaf 9 bezoeken per maand is Sterk goedkoper.} Controle: bij 9 bezoeken kost Sterk 43 euro en Fit 45 euro, bij 8 bezoeken 41 en 40 euro.}

In de figuur zie je het ook: tot 8 bezoeken liggen de vierkantjes van Fit onder de bolletjes van Sterk, vanaf 9 erboven. De ongelijkheid gaf $x > 8{,}3$, maar het antwoord moet een geheel aantal zijn. Dan moet je afronden, en de richting hangt af van de vraag.

\begin{werkwijze}
Kies de onbekende en schrijf erbij wat ze betekent. Zoek in de tekst het woord dat het teken bepaalt (minstens, hoogstens, goedkoper, \ldots) en stel de ongelijkheid op. Los ze op. Kijk welke waarden in de context kunnen: moet het een geheel getal zijn, of positief? Rond in de goede richting af, controleer, en antwoord met een zin.
\end{werkwijze}

\voorbeeld{}{Een lift mag hoogstens 630~kg dragen. Een verhuizer van 85~kg wil dozen van 24~kg meenemen. Hoeveel dozen mag hij hoogstens meenemen?\label{vb:h2-lift}}{%
\uitleg{\textit{Hoogstens} geeft het teken $\leqslant$. Het totale gewicht is de verhuizer plus de dozen.}
\begin{stappen}
x = \text{het aantal dozen} & De onbekende. \\
85 + 24x \leqslant 630 & Het totale gewicht is hoogstens 630~kg. \\
24x \leqslant 545 \ \Rightarrow \ x \leqslant 22{,}70\ldots & Breng 85 over en deel door 24.
\end{stappen}
Het aantal dozen is een geheel getal, en hoogstens $22{,}70\ldots$: \textit{hij mag hoogstens 22 dozen meenemen.} Hier rond je naar beneden af, ook al ligt $22{,}7$ dichter bij 23. Controle: $85 + 22 \cdot 24 = 613 \leqslant 630$, maar met 23 dozen wordt het 637~kg.}

\voorbeeld{}{Een klas verkoopt wafels voor een goed doel. Ze huurt een wafelijzer voor 35 euro. Elke wafel kost 0,60 euro aan ingrediënten en wordt verkocht voor 2,50 euro. Hoeveel wafels moet de klas minstens verkopen om minstens 400 euro winst te maken?\label{vb:h2-wafels}}{%
\uitleg{De winst is de opbrengst min alle kosten. \textit{Minstens} geeft $\geqslant$.}
\begin{stappen}
x = \text{het aantal wafels} & De onbekende. \\
2{,}5x - 0{,}6x - 35 \geqslant 400 & Opbrengst min ingrediënten min huur. \\
1{,}9x \geqslant 435 \ \Rightarrow \ x \geqslant 228{,}94\ldots & Neem samen, breng over en deel door $1{,}9$.
\end{stappen}
\textit{De klas moet minstens 229 wafels verkopen.} Nu rond je naar boven af: met 228 wafels is de winst $398{,}20$ euro, met 229 wafels $400{,}10$ euro.}

\begin{oefeningen}
\oef{h2-vraagstuk}
\end{oefeningen}

%========================================================================================
\opdrachtenkop

\begin{oefeningen}
\oef{h2-o1}

\oef{h2-o2}

\oef{h2-o3}

\oef{h2-o4}

\oef{h2-o5}

\oef{h2-o6}

\oef{h2-o7}

\oef{h2-o8}

\oef{h2-o9}

\oef{h2-o10}

\oef{h2-o11}

\oef{h2-o12}

\oef{h2-o13}

\oef{h2-o14}

\oef{h2-o15}

\oef{h2-o16}

\oef{h2-o17}

\oef{h2-o18}

\oef{h2-o19}

\oef{h2-o20}

\oef{h2-o21}

\oef{h2-o22}

\oef{h2-o23}

\oef{h2-o24}
\end{oefeningen}

%----------------------------------------------------------------------------------------
\needspace{14\baselineskip}
\opdrachtsectie{Test jezelf}

\begin{oefeningen}
\oef{h2-t1}

\oef{h2-t2}

\oef{h2-t3}

\oef{h2-t4}

\oef{h2-t5}

\oef{h2-t6}

\oef{h2-t7}

\oef{h2-t8}
\end{oefeningen}

\herhaalschema{%
  \fv{oef:h2-t1}{vb:h2-interval}{oef:h2-interval}%
  \fv{oef:h2-t2}{vb:h2-positief}{oef:h2-positief}%
  \fv{oef:h2-t3}{vb:h2-negatief}{oef:h2-negatief}%
  \fv{oef:h2-t4}{vb:h2-herleid}{oef:h2-herleid}%
  \fv{oef:h2-t5}{vb:h2-geen}{oef:h2-o11}%
  \fv{oef:h2-t6}{vb:h2-lift}{oef:h2-o17}%
  \fv{oef:h2-t7}{vb:h2-terug}{oef:h2-o14}%
  \alleenvoor{4,5}{\fv{oef:h2-t8}{vb:h2-dubbel}{oef:h2-dubbel}}}

\ifzelfstudie
\section*{Maak je eigen samenvatting}
Sluit dit boek en schrijf in je schrift: wat een ongelijkheid en een oplossing zijn, de soorten intervallen met hun notatie en tekening, de drie rekenregels (en wanneer het teken omkeert), \alleenvoor{4,5}{hoe je een dubbele ongelijkheid oplost, }en hoe je bij een vraagstuk in de goede richting afrondt. Vergelijk daarna met de kaders.
\fi
